\chapter{Quiz: scelta multipla e vero/falso}
\label{cap:quiz}
\setcounter{domanda}{0}

Domande nello stile delle prime due parti dello scritto (cap.~\ref{cap:formato}): scelta multipla e
gruppi di vero/falso. \textbf{Le soluzioni sono in fondo al capitolo} (sez.~\ref{sec:quiz-soluzioni}),
ciascuna con il rimando alla lezione in cui l'argomento è trattato: prova prima da solo.

\begin{nota}[Domande dei redattori]
Sono domande scritte per il ripasso a partire da lezioni e slide, \textbf{non} domande ufficiali
d'esame. Il formato reale si vedrà nella simulazione d'esame dell'ultima settimana.
\end{nota}

\section{Domande a scelta multipla}
Una sola risposta corretta per domanda.

\begin{domanda}[AI e ML]\label{quiz:m1}
Quale affermazione su AI e machine learning è corretta?
\begin{opzioni}
  \item AI e machine learning sono sinonimi.
  \item Il machine learning non richiede dati: la regola la fornisce l'esperto.
  \item Un sistema a regole scritte da un esperto è AI ma non machine learning.
  \item Il deep learning è un sottoinsieme dell'AI che non usa reti neurali.
\end{opzioni}
\end{domanda}

\begin{domanda}[shallow vs deep]\label{quiz:m2}
Cosa distingue principalmente il deep learning dallo shallow learning, secondo il corso?
\begin{opzioni}
  \item La feature extraction è fatta dalla rete; lo sforzo si sposta sul model engineering.
  \item Il deep learning non richiede preprocessing.
  \item Il deep learning è più interpretabile.
  \item Il deep learning si applica solo alle immagini.
\end{opzioni}
\end{domanda}

\begin{domanda}[terminologia]\label{quiz:m3}
Il validation set serve a:
\begin{opzioni}
  \item aggiornare i pesi della rete;
  \item aumentare la quantità di dati di training;
  \item stimare le prestazioni finali dopo aver fissato tutti gli iperparametri;
  \item scegliere gli iperparametri e valutare l'overfitting.
\end{opzioni}
\end{domanda}

\begin{domanda}[paradigmi]\label{quiz:m4}
Qual è la differenza tra unsupervised e self-supervised learning?
\begin{opzioni}
  \item Il self-supervised richiede label umane, l'unsupervised no.
  \item Il self-supervised genera un segnale di supervisione dai dati stessi per imparare una rappresentazione.
  \item L'unsupervised impara una rappresentazione, il self-supervised descrive i dati.
  \item Sono sinonimi.
\end{opzioni}
\end{domanda}

\begin{domanda}[federated learning]\label{quiz:m5}
Nel federated learning, il \en{model poisoning} consiste nel:
\begin{opzioni}
  \item iniettare rumore impercettibile nelle immagini di test;
  \item ricostruire i dati di training a partire dal modello;
  \item far inviare a un partecipante un modello addestrato per far passare come legittimo il proprio traffico malevolo;
  \item aggirare i guardrail di un LLM con un prompt.
\end{opzioni}
\end{domanda}

\begin{domanda}[adversarial ML]\label{quiz:m6}
Un \en{adversarial example}:
\begin{opzioni}
  \item modifica leggermente l'input perché il modello sbagli, senza toccare il modello;
  \item modifica i pesi del modello durante il training;
  \item ricostruisce i dati di training dagli output del modello;
  \item è un campione con la label sbagliata.
\end{opzioni}
\end{domanda}

\begin{domanda}[perceptron]\label{quiz:m7}
A cosa serve il bias $w_0$ in un perceptron?
\begin{opzioni}
  \item A rendere la funzione non lineare.
  \item A normalizzare l'input.
  \item A ridurre il numero di parametri.
  \item A traslare la funzione, che altrimenti passerebbe sempre per l'origine.
\end{opzioni}
\end{domanda}

\begin{domanda}[non linearità]\label{quiz:m8}
Una rete con 10 layer lineari e nessuna funzione di attivazione:
\begin{opzioni}
  \item può approssimare qualsiasi funzione continua;
  \item equivale a un unico layer lineare;
  \item è più espressiva di una rete con un solo hidden layer ReLU;
  \item soffre di over-smoothing.
\end{opzioni}
\end{domanda}

\begin{domanda}[output layer]\label{quiz:m9}
Classificazione con 5 classi: quale output layer?
\begin{opzioni}
  \item 1 neurone con sigmoid.
  \item 5 neuroni con ReLU.
  \item 5 neuroni con softmax.
  \item 1 neurone con tanh.
\end{opzioni}
\end{domanda}

\begin{domanda}[ottimizzazione]\label{quiz:m10}
Rispetto al gradient descent, lo stochastic gradient descent:
\begin{opzioni}
  \item aggiorna i pesi dopo ogni mini-batch e, grazie al rumore, può uscire dai minimi locali;
  \item usa tutto il training set per ogni aggiornamento;
  \item garantisce di trovare il minimo globale;
  \item non richiede un learning rate.
\end{opzioni}
\end{domanda}

\begin{domanda}[momentum]\label{quiz:m11}
Con un coefficiente di momentum $\beta$ molto vicino a 1:
\begin{opzioni}
  \item l'aggiornamento segue quasi solo il gradiente del batch corrente;
  \item il learning rate effettivo diventa nullo;
  \item l'ottimizzatore equivale all'SGD senza momentum;
  \item l'aggiornamento segue soprattutto la direzione passata: traiettoria liscia ma lenta a cambiare direzione.
\end{opzioni}
\end{domanda}

\begin{domanda}[Adam]\label{quiz:m12}
La \en{bias correction} di Adam serve a:
\begin{opzioni}
  \item correggere i bias dei neuroni;
  \item compensare il fatto che $m$ e $v$ partono da 0 e sono sottostimati nei primi passi;
  \item aumentare il learning rate nel tempo;
  \item prevenire l'overfitting.
\end{opzioni}
\end{domanda}

\begin{domanda}[vanishing gradient]\label{quiz:m13}
Perché la sigmoid negli hidden layer di una rete profonda è problematica?
\begin{opzioni}
  \item Non è derivabile.
  \item Non è limitata e fa esplodere il gradiente.
  \item Le sue derivate sono piccole e, moltiplicate layer dopo layer, fanno svanire il gradiente.
  \item Aggiunge troppi parametri.
\end{opzioni}
\end{domanda}

\begin{domanda}[metriche]\label{quiz:m14}
Dataset con 9.800 campioni benigni e 200 malware. Un classificatore che risponde sempre ``benigno'' ha:
\begin{opzioni}
  \item accuracy 98\% e recall sul malware pari a 0;
  \item accuracy 2\% e recall sul malware pari a 1;
  \item precision sul malware pari a 1;
  \item F1 sul malware pari a 0,98.
\end{opzioni}
\end{domanda}

\begin{domanda}[learning curve]\label{quiz:m15}
Learning curve: la loss di training scende quasi a 0, quella di validation scende fino all'epoch 30 e poi risale. Cosa fai?
\begin{opzioni}
  \item Aumento il numero di neuroni.
  \item Aumento il learning rate.
  \item Continuo il training finché anche la validation loss arriva a 0.
  \item Uso l'early stopping e ripristino i pesi salvati attorno all'epoch 30.
\end{opzioni}
\end{domanda}

\begin{domanda}[regolarizzazione]\label{quiz:m16}
La regolarizzazione L2:
\begin{opzioni}
  \item aggiunge alla loss $\lambda\sum_i|\phi_i|$, spingendo molti pesi a 0;
  \item aggiunge alla loss $\lambda\sum_i\phi_i^2$, scoraggiando i pesi grandi;
  \item spegne neuroni a caso durante il training;
  \item normalizza gli input di ogni layer.
\end{opzioni}
\end{domanda}

\begin{domanda}[scaling]\label{quiz:m17}
Quando la standardizzazione (z-score) è preferibile alla normalizzazione min-max?
\begin{opzioni}
  \item Quando servono valori nell'intervallo $[0,1]$.
  \item Per le feature categoriche.
  \item Quando ci sono outlier.
  \item Solo per la label.
\end{opzioni}
\end{domanda}

\begin{domanda}[scaling e leakage]\label{quiz:m18}
Su quali dati si esegue il \code{fit} dello \code{StandardScaler}?
\begin{opzioni}
  \item Solo sul training set.
  \item Su tutto il dataset.
  \item Su training e validation.
  \item Sul test set.
\end{opzioni}
\end{domanda}

\begin{domanda}[feature categoriche]\label{quiz:m19}
La porta TCP di destinazione di un flusso è meglio trattata come feature:
\begin{opzioni}
  \item numerica standardizzata;
  \item numerica normalizzata min-max;
  \item label;
  \item categorica.
\end{opzioni}
\end{domanda}

\begin{domanda}[cross-validation]\label{quiz:m20}
Con la k-fold cross-validation, una volta scelti gli iperparametri, il modello finale:
\begin{opzioni}
  \item è quello del fold con la validation migliore;
  \item si ri-addestra su tutti i fold con gli iperparametri scelti;
  \item è la media dei pesi dei $k$ modelli;
  \item si addestra sul test set.
\end{opzioni}
\end{domanda}

\begin{domanda}[data leakage]\label{quiz:m21}
Quale di queste pratiche introduce data leakage?
\begin{opzioni}
  \item Stratificare lo split.
  \item Usare il dropout.
  \item Includere tra le feature un identificativo univoco del campione.
  \item Fare early stopping sulla validation loss.
\end{opzioni}
\end{domanda}

\begin{domanda}[PyTorch: loss]\label{quiz:m22}
Con \code{nn.CrossEntropyLoss}, l'ultimo layer del modello deve restituire:
\begin{opzioni}
  \item i logit grezzi;
  \item le probabilità dopo la softmax;
  \item vettori one-hot;
  \item le probabilità dopo la sigmoid.
\end{opzioni}
\end{domanda}

\begin{domanda}[PyTorch: training loop]\label{quiz:m23}
In un training loop con mini-batch, \code{optimizer.zero\_grad()}, \code{loss.backward()} e \code{optimizer.step()} vanno:
\begin{opzioni}
  \item fuori dal ciclo sui batch, una volta per epoch;
  \item solo nel ciclo di validation;
  \item prima del ciclo sulle epoch;
  \item dentro il ciclo sui mini-batch.
\end{opzioni}
\end{domanda}

\begin{domanda}[PyTorch: eval]\label{quiz:m24}
Cosa fa \code{model.eval()}?
\begin{opzioni}
  \item Disattiva il calcolo dei gradienti.
  \item Mette dropout e batch norm in modalità di valutazione.
  \item Congela definitivamente i pesi.
  \item Calcola accuracy e F1.
\end{opzioni}
\end{domanda}

\begin{domanda}[RNN]\label{quiz:m25}
Classificare un'email (disponibile per intero) come phishing a partire dalla sequenza delle sue parole è un caso:
\begin{opzioni}
  \item one-to-many, solo con RNN mono-direzionale;
  \item one-to-one;
  \item many-to-one, in cui si può usare una RNN bidirezionale;
  \item many-to-many in tempo reale.
\end{opzioni}
\end{domanda}

\begin{domanda}[LSTM]\label{quiz:m26}
Nell'LSTM, il \en{forget gate}:
\begin{opzioni}
  \item decide con una sigmoid quanto del cell state precedente mantenere;
  \item aggiunge nuova informazione al cell state con una tanh;
  \item produce direttamente l'output della rete;
  \item azzera l'hidden state a ogni passo.
\end{opzioni}
\end{domanda}

\begin{domanda}[tensori RNN]\label{quiz:m27}
Un batch di 16 serie univariate lunghe 50, con \code{batch\_first=True}, deve avere forma:
\begin{opzioni}
  \item \code{(16, 50)};
  \item \code{(50, 16)};
  \item \code{(16, 1, 50)};
  \item \code{(16, 50, 1)}.
\end{opzioni}
\end{domanda}

\begin{domanda}[CNN]\label{quiz:m28}
In una CNN, il max pooling:
\begin{opzioni}
  \item aumenta la dimensione spaziale;
  \item riduce la dimensione sostituendo ogni finestra con il suo massimo;
  \item aggiunge pesi addestrabili;
  \item sostituisce la ReLU come non linearità.
\end{opzioni}
\end{domanda}

\begin{domanda}[GNN]\label{quiz:m29}
Perché una FFNN applicata alla matrice di adiacenza non è adatta ai grafi?
\begin{opzioni}
  \item Perché i grafi non hanno feature.
  \item Perché è troppo lenta da addestrare.
  \item Perché la rappresentazione dipende dall'ordine dei nodi e non gestisce grafi di dimensioni diverse.
  \item Perché non usa funzioni di attivazione.
\end{opzioni}
\end{domanda}

\begin{domanda}[GNN]\label{quiz:m30}
L'over-smoothing in una GNN è:
\begin{opzioni}
  \item l'appiattimento degli embedding dei nodi, che diventano simili perché i receptive field si sovrappongono;
  \item una loss di training che oscilla;
  \item un'esplosione dei gradienti;
  \item la perdita di memoria del passato lontano.
\end{opzioni}
\end{domanda}

\begin{domanda}[GNN]\label{quiz:m31}
GraphSAGE si distingue perché:
\begin{opzioni}
  \item usa coefficienti di attenzione appresi;
  \item usa tutta la matrice di adiacenza normalizzata per il grado;
  \item non condivide i pesi tra i nodi;
  \item campiona un numero fisso di vicini.
\end{opzioni}
\end{domanda}

\begin{domanda}[PyTorch Geometric]\label{quiz:m32}
Nel training di una GNN per graph classification, la loss del batch va pesata per:
\begin{opzioni}
  \item \code{batch.x.size(0)};
  \item \code{batch.num\_graphs};
  \item il numero di archi;
  \item la hidden size.
\end{opzioni}
\end{domanda}

\begin{domanda}[embedding]\label{quiz:m33}
Per una feature nominale ad alta cardinalità (es.\ il dominio del mittente) in una rete neurale, la scelta migliore è:
\begin{opzioni}
  \item label encoding in ordine alfabetico;
  \item one-hot encoding;
  \item un embedding appreso, con vocabolario costruito sul training e token \code{<UNK>};
  \item eliminarla.
\end{opzioni}
\end{domanda}

\begin{domanda}[embedding]\label{quiz:m34}
\code{nn.Embedding(V, H)}:
\begin{opzioni}
  \item crea una lookup addestrabile con $V$ righe e $H$ colonne, che riceve indici interi;
  \item crea una matrice non addestrabile con $H$ righe e $V$ colonne;
  \item riceve in input le categorie come numeri float;
  \item si può usare anche con un random forest.
\end{opzioni}
\end{domanda}

\begin{domanda}[dropout]\label{quiz:m35}
Dropout con $p = 0{,}5$ in training: un neurone che sopravvive con output 3 produce:
\begin{opzioni}
  \item 3;
  \item 1,5;
  \item 0;
  \item 6.
\end{opzioni}
\end{domanda}

\begin{domanda}[batch normalization]\label{quiz:m36}
In validation, la batch normalization usa:
\begin{opzioni}
  \item media e varianza del batch di validation;
  \item running mean e running variance memorizzate durante il training;
  \item nessuna normalizzazione;
  \item media e varianza del test set.
\end{opzioni}
\end{domanda}

\begin{domanda}[tipi di anomalia]\label{quiz:m37}
Backup di 2 ore la domenica, quando l'azienda è chiusa; nei giorni lavorativi backup di 2 ore sono normali. È un'anomalia:
\begin{opzioni}
  \item puntuale;
  \item collettiva;
  \item contestuale;
  \item non è un'anomalia.
\end{opzioni}
\end{domanda}

\begin{domanda}[one-class SVM]\label{quiz:m38}
Nella one-class SVM non supervisionata, $\nu$ rappresenta:
\begin{opzioni}
  \item la frazione attesa di outlier;
  \item il numero di cluster;
  \item il grado del kernel polinomiale;
  \item il learning rate.
\end{opzioni}
\end{domanda}

\begin{domanda}[k-means]\label{quiz:m39}
Nel grafico SSE in funzione di $K$ (k-means), l'SSE risale passando da $K=5$ a $K=6$. Causa probabile?
\begin{opzioni}
  \item $K=6$ è il valore ottimo.
  \item I dati non hanno struttura a cluster.
  \item La distanza euclidea non è adatta.
  \item Un'inizializzazione sfortunata (es.\ un cluster vuoto): va rieseguito k-means.
\end{opzioni}
\end{domanda}

\begin{domanda}[anomaly detection]\label{quiz:m40}
Vantaggio di un algoritmo di anomaly detection che restituisce \textbf{score} invece di label:
\begin{opzioni}
  \item non richiede dati di training;
  \item permette di ordinare le istanze e scegliere un budget di analisi;
  \item è sempre più accurato;
  \item elimina i falsi positivi.
\end{opzioni}
\end{domanda}

\begin{domanda}[PCA]\label{quiz:m41}
Su dati non scalati, la PCA dice che una sola componente spiega il 99,9\% della varianza. Causa più probabile?
\begin{opzioni}
  \item Il dataset è davvero unidimensionale.
  \item Ci sono troppe feature categoriche in one-hot.
  \item Una feature numerica con scala enorme (es.\ byte) domina la varianza: manca lo scaling.
  \item La PCA andava sostituita con t-SNE.
\end{opzioni}
\end{domanda}

\begin{domanda}[autoencoder]\label{quiz:m42}
Perché un autoencoder addestrato solo su dati normali rileva le anomalie meglio di uno addestrato su normali e anomalie?
\begin{opzioni}
  \item Con le anomalie in training impara a ricostruire anche quelle e il loro errore di ricostruzione si avvicina a quello dei normali.
  \item Con meno dati il training converge sempre a un minimo migliore.
  \item Solo con i dati normali si può usare la cross-entropy.
  \item Le anomalie in training causano sempre vanishing gradient.
\end{opzioni}
\end{domanda}

\begin{domanda}[Isolation Forest]\label{quiz:m43}
Quale di questi metodi non si basa su distanze o densità, e soffre quindi meno l'alta dimensionalità?
\begin{opzioni}
  \item LOF.
  \item DBSCAN.
  \item K-means con distanza dal centroide.
  \item Isolation Forest.
\end{opzioni}
\end{domanda}

\begin{domanda}[LOF]\label{quiz:m44}
Un attacco forma un gruppo piccolo ma molto denso di connessioni quasi identiche. Con $k=10$ i suoi punti hanno LOF $\approx 1$. Significa che:
\begin{opzioni}
  \item i punti sono outlier certi;
  \item hanno una densità simile ai vicini, quindi LOF non li segnala;
  \item il modello non è stato addestrato;
  \item la contaminazione è troppo alta.
\end{opzioni}
\end{domanda}


\begin{domanda}[bag of words]\label{quiz:m45}
Quale affermazione sulla bag of words (count vectorizer) è corretta?
\begin{opzioni}
  \item Conserva l'ordine delle parole ma non la loro frequenza.
  \item Rappresenta ogni frase con i conteggi delle parole del vocabolario: perde ordine e semantica, e la dimensione cresce con il vocabolario.
  \item Richiede una rete neurale per essere calcolata.
  \item Assegna a parole con significato simile vettori vicini.
\end{opzioni}
\end{domanda}

\begin{domanda}[Word2Vec]\label{quiz:m46}
Nello skip-gram di Word2Vec la rete impara a:
\begin{opzioni}
  \item predire la parola centrale a partire dalle parole del contesto;
  \item classificare le frasi come positive o negative;
  \item predire le parole del contesto a partire dalla parola centrale;
  \item ricostruire la frase intera a partire dal suo embedding.
\end{opzioni}
\end{domanda}

\begin{domanda}[Word2Vec]\label{quiz:m47}
Nella probabilità $p(w_{t+j} \mid w_t) = \exp(e_{t+j}^{\top} e_t) / \sum_{w=1}^{V} \exp(e_w^{\top} e_t)$, cosa rende il calcolo impraticabile e come lo si risolve?
\begin{opzioni}
  \item Il numeratore, perché il prodotto scalare è costoso; lo si sostituisce con la distanza euclidea.
  \item Il denominatore, che somma su tutto il vocabolario a ogni passo; lo si approssima con $k$ parole negative estratte a caso.
  \item Il contesto $c$, che deve coprire tutta la frase; lo si riduce a una parola.
  \item L'esponenziale, che va in overflow; lo si toglie.
\end{opzioni}
\end{domanda}

\begin{domanda}[Word2Vec]\label{quiz:m48}
Un corpus contiene in parti uguali testi di zoologia, favole e libri di informatica. Che embedding avrà la parola \emph{mouse} in Word2Vec?
\begin{opzioni}
  \item Due embedding, uno per l'animale e uno per il dispositivo.
  \item Un embedding diverso a seconda della frase in cui la parola compare.
  \item Un solo embedding, spostato verso il significato più frequente (l'animale), che non rappresenta bene nessuno dei due.
  \item Nessun embedding, perché la parola è ambigua.
\end{opzioni}
\end{domanda}

\begin{domanda}[darknet]\label{quiz:m49}
Perché tutto il traffico ricevuto da una darknet è interessante?
\begin{opzioni}
  \item Perché contiene il traffico degli utenti legittimi verso i servizi dell'ateneo.
  \item Perché la darknet risponde agli scanner e ne registra le richieste complete.
  \item Perché i suoi IP non offrono servizi né inviano richieste: ogni pacchetto è non richiesto (scansioni, backscattering, errori).
  \item Perché ogni pacchetto arriva con una label che indica l'attaccante.
\end{opzioni}
\end{domanda}

\begin{domanda}[DarkVec]\label{quiz:m50}
In DarkVec, che cosa corrisponde alle ``parole'' e alle ``frasi'' di Word2Vec?
\begin{opzioni}
  \item Parole = porte di destinazione; frasi = pacchetti di un singolo IP.
  \item Parole = IP sorgente; frasi = sequenze di IP sorgente in ordine di arrivo, una per servizio (porta/protocollo).
  \item Parole = byte del payload; frasi = sessioni TCP complete.
  \item Parole = IP di destinazione della darknet; frasi = giorni di osservazione.
\end{opzioni}
\end{domanda}

\begin{domanda}[i-DarkVec]\label{quiz:m51}
Quale limite di Word2Vec affronta i-DarkVec?
\begin{opzioni}
  \item L'impossibilità di usare il negative sampling.
  \item Gli IP nuovi (fuori vocabolario) e i comportamenti che cambiano, aggiornando gli embedding giorno per giorno a partire da quelli già noti.
  \item La necessità di label per addestrare gli embedding.
  \item La dimensione eccessiva degli embedding rispetto al one-hot.
\end{opzioni}
\end{domanda}

\section{Vero o falso}
Come all'esame: risposta giusta = punto, risposta sbagliata = piccola penalità. Se non sei sicuro,
valuta se rischiare. Per ogni affermazione indica V o F.

\begin{domanda}[vero/falso -- Training]\label{quiz:vfA}
Gruppo A -- Training.
\begin{enumerate}[label=A\arabic*.,nosep,leftmargin=2.6em]
  \item Un'epoch corrisponde a un singolo aggiornamento dei pesi anche quando si usano i mini-batch. \hfill V\quad F
  \item I mini-batch si campionano senza reinserimento, così ogni punto è usato una sola volta per epoch. \hfill V\quad F
  \item Una loss di training pari a 0 indica che il modello è pronto per la produzione. \hfill V\quad F
  \item Un learning rate scheduler parte tipicamente da valori alti e li riduce durante il training. \hfill V\quad F
\end{enumerate}
\end{domanda}

\begin{domanda}[vero/falso -- Dati]\label{quiz:vfB}
Gruppo B -- Dati.
\begin{enumerate}[label=B\arabic*.,nosep,leftmargin=2.6em]
  \item Il label encoding è adatto alla label di classificazione perché la cross-entropy non usa distanze tra classi. \hfill V\quad F
  \item Il one-hot encoding introduce un ordine tra le categorie. \hfill V\quad F
  \item Con dati temporali, uno split casuale può causare data leakage. \hfill V\quad F
  \item Con classi molto sbilanciate, la media \en{micro} delle metriche dà lo stesso peso a ogni classe. \hfill V\quad F
\end{enumerate}
\end{domanda}

\begin{domanda}[vero/falso -- PyTorch]\label{quiz:vfC}
Gruppo C -- PyTorch.
\begin{enumerate}[label=C\arabic*.,nosep,leftmargin=2.6em]
  \item \code{torch.no\_grad()} spegne il dropout. \hfill V\quad F
  \item Nel DataLoader di validation lo shuffle non serve. \hfill V\quad F
  \item Se il modello è sulla GPU, anche i tensori dei batch devono essere spostati sullo stesso device. \hfill V\quad F
  \item Con \code{nn.CrossEntropyLoss} le label devono essere necessariamente vettori one-hot di tipo float. \hfill V\quad F
\end{enumerate}
\end{domanda}

\begin{domanda}[vero/falso -- RNN e CNN]\label{quiz:vfD}
Gruppo D -- RNN e CNN.
\begin{enumerate}[label=D\arabic*.,nosep,leftmargin=2.6em]
  \item In una RNN la matrice $W_{hh}$ è diversa per ogni passo temporale. \hfill V\quad F
  \item Una RNN bidirezionale è adatta alla predizione in tempo reale. \hfill V\quad F
  \item Il gradient clipping è un rimedio all'exploding gradient. \hfill V\quad F
  \item In una CNN i valori dei filtri convoluzionali sono pesi appresi durante il training. \hfill V\quad F
\end{enumerate}
\end{domanda}

\begin{domanda}[vero/falso -- GNN ed embedding]\label{quiz:vfE}
Gruppo E -- GNN ed embedding.
\begin{enumerate}[label=E\arabic*.,nosep,leftmargin=2.6em]
  \item Una GNN con aggregazione a somma è invariante all'ordine dei vicini. \hfill V\quad F
  \item Le GNN non possono produrre embedding per nodi mai visti in training. \hfill V\quad F
  \item Il vocabolario per gli embedding va costruito su training e validation insieme, per ridurre i \code{<UNK>}. \hfill V\quad F
  \item Gli embedding appresi con supervisione sono ottimi per il task di training ma possono essere sub-ottimi per altri task. \hfill V\quad F
\end{enumerate}
\end{domanda}

\begin{domanda}[vero/falso -- Anomaly detection]\label{quiz:vfF}
Gruppo F -- Anomaly detection.
\begin{enumerate}[label=F\arabic*.,nosep,leftmargin=2.6em]
  \item Una novità può diventare la nuova normalità. \hfill V\quad F
  \item Gli algoritmi shallow riconoscono le anomalie collettive senza bisogno di feature engineering. \hfill V\quad F
  \item K-means assegna ogni punto a un cluster, quindi non produce direttamente ``punti di rumore''. \hfill V\quad F
  \item Nella one-class SVM ``supervisionata'' il modello riceve le label normale/anomalo. \hfill V\quad F
\end{enumerate}
\end{domanda}

\begin{domanda}[vero/falso -- Anomaly detection in pratica]\label{quiz:vfG}
Gruppo G -- Anomaly detection in pratica.
\begin{enumerate}[label=G\arabic*.,nosep,leftmargin=2.6em]
  \item Nell'anomaly detection le label disponibili si usano per addestrare il modello. \hfill V\quad F
  \item Nella mappa di t-SNE la distanza tra due gruppi misura quanto sono diversi nello spazio originale. \hfill V\quad F
  \item DBSCAN etichetta come rumore i punti che non sono né core né border point. \hfill V\quad F
  \item Su dati molto sbilanciati una ROC-AUC alta garantisce un buon F1 sulla classe delle anomalie. \hfill V\quad F
\end{enumerate}
\end{domanda}


\begin{domanda}[vero/falso -- NLP e DarkVec]\label{quiz:vfH}
Gruppo H -- NLP e DarkVec.
\begin{enumerate}[label=H\arabic*.,nosep,leftmargin=2.6em]
  \item Word2Vec ha bisogno di frasi etichettate a mano per imparare gli embedding. \hfill V\quad F
  \item Una finestra di contesto $c$ più grande porta più informazione ma richiede più calcolo. \hfill V\quad F
  \item Word2Vec produce embedding contestuali: la stessa parola ha vettori diversi in frasi diverse. \hfill V\quad F
  \item In DarkVec l'ipotesi è che gli IP che compaiono insieme nel tempo sugli stessi servizi abbiano un comportamento simile. \hfill V\quad F
\end{enumerate}
\end{domanda}

\section{Soluzioni}
\label{sec:quiz-soluzioni}


\begin{soluzione}[Domanda~\ref{quiz:m1}]
\textbf{(c)}. L'AI include i sistemi a regole; il ML impara dai dati. (rif.\ \rif{2026-09-21}{21/09})
\end{soluzione}

\begin{soluzione}[Domanda~\ref{quiz:m2}]
\textbf{(a)}. Il preprocessing resta; i modelli deep sono meno interpretabili. (rif.\ \rif{2026-09-21}{21/09})
\end{soluzione}

\begin{soluzione}[Domanda~\ref{quiz:m3}]
\textbf{(d)}. (c) è il ruolo del test set. (rif.\ \rif{2026-09-23}{23/09})
\end{soluzione}

\begin{soluzione}[Domanda~\ref{quiz:m4}]
\textbf{(b)}. Unsupervised = descrizione; self-supervised = rappresentazione con supervisione dai dati. (rif.\ \rif{2026-09-23}{23/09})
\end{soluzione}

\begin{soluzione}[Domanda~\ref{quiz:m5}]
\textbf{(c)}. (b) è model inversion, (d) jailbreak, (a) adversarial example. (rif.\ \rif{2026-09-23}{23/09})
\end{soluzione}

\begin{soluzione}[Domanda~\ref{quiz:m6}]
\textbf{(a)}. (rif.\ \rif{2026-09-23}{23/09})
\end{soluzione}

\begin{soluzione}[Domanda~\ref{quiz:m7}]
\textbf{(d)}. Senza bias, con attivazione lineare, ogni retta passa per l'origine. (rif.\ \rif{2026-09-28}{28/09})
\end{soluzione}

\begin{soluzione}[Domanda~\ref{quiz:m8}]
\textbf{(b)}. Prodotto di matrici = una matrice: il modello resta lineare. (rif.\ \rif{2026-09-28}{28/09})
\end{soluzione}

\begin{soluzione}[Domanda~\ref{quiz:m9}]
\textbf{(c)}. (rif.\ \rif{2026-09-28}{28/09})
\end{soluzione}

\begin{soluzione}[Domanda~\ref{quiz:m10}]
\textbf{(a)}. L'SGD non garantisce il minimo globale: rende solo più probabile uscire da quelli locali. (rif.\ \rif{2026-09-30}{30/09})
\end{soluzione}

\begin{soluzione}[Domanda~\ref{quiz:m11}]
\textbf{(d)}. $m_{t+1} = \beta m_t + (1-\beta)g$: con $\beta\to1$ il gradiente nuovo pesa pochissimo. (rif.\ \rif{2026-09-30}{30/09})
\end{soluzione}

\begin{soluzione}[Domanda~\ref{quiz:m12}]
\textbf{(b)}. (rif.\ \rif{2026-09-30}{30/09})
\end{soluzione}

\begin{soluzione}[Domanda~\ref{quiz:m13}]
\textbf{(c)}. La derivata della sigmoid vale al massimo 0,25. (rif.\ \rif{2026-09-30}{30/09})
\end{soluzione}

\begin{soluzione}[Domanda~\ref{quiz:m14}]
\textbf{(a)}. $9800/10000 = 98\%$; nessun malware trovato $\Rightarrow$ recall 0; precision 0/0 non definita. (rif.\ \rif{2026-09-30}{30/09})
\end{soluzione}

\begin{soluzione}[Domanda~\ref{quiz:m15}]
\textbf{(d)}. Il modello dell'ultima epoch è già in overfitting: serve lo stato salvato. (rif.\ \rif{2026-09-30}{30/09})
\end{soluzione}

\begin{soluzione}[Domanda~\ref{quiz:m16}]
\textbf{(b)}. (a) è L1; (c) dropout; (d) batch norm. (rif.\ \rif{2026-09-30}{30/09})
\end{soluzione}

\begin{soluzione}[Domanda~\ref{quiz:m17}]
\textbf{(c)}. (rif.\ \rif{2026-10-05}{05/10})
\end{soluzione}

\begin{soluzione}[Domanda~\ref{quiz:m18}]
\textbf{(a)}. (rif.\ \rif{2026-10-05}{05/10})
\end{soluzione}

\begin{soluzione}[Domanda~\ref{quiz:m19}]
\textbf{(d)}. 80, 81, 8080: la distanza numerica non ha significato. (rif.\ \rif{2026-10-05}{05/10})
\end{soluzione}

\begin{soluzione}[Domanda~\ref{quiz:m20}]
\textbf{(b)}. Scegliere il fold migliore significa sfruttare la fortuna sul fold di validation. (rif.\ \rif{2026-10-05}{05/10})
\end{soluzione}

\begin{soluzione}[Domanda~\ref{quiz:m21}]
\textbf{(c)}. Un ID permette di memorizzare i campioni; lo split temporale casuale è un altro esempio. (rif.\ \rif{2026-10-05}{05/10})
\end{soluzione}

\begin{soluzione}[Domanda~\ref{quiz:m22}]
\textbf{(a)}. La softmax è dentro la loss; applicarla due volte appiattisce i gradienti. (rif.\ \rif{2026-10-06}{06/10})
\end{soluzione}

\begin{soluzione}[Domanda~\ref{quiz:m23}]
\textbf{(d)}. Fuori dal ciclo si farebbe un solo aggiornamento per epoch (GD con gradienti sommati). (rif.\ \rif{2026-10-06}{06/10})
\end{soluzione}

\begin{soluzione}[Domanda~\ref{quiz:m24}]
\textbf{(b)}. \code{no\_grad} spegne i gradienti; \code{eval} cambia il comportamento di dropout e BN. (rif.\ \rif{2026-10-06}{06/10})
\end{soluzione}

\begin{soluzione}[Domanda~\ref{quiz:m25}]
\textbf{(c)}. L'email è completa: si può leggere anche all'indietro; un output per tutta la sequenza. (rif.\ \rif{2026-10-12}{12/10})
\end{soluzione}

\begin{soluzione}[Domanda~\ref{quiz:m26}]
\textbf{(a)}. (rif.\ \rif{2026-10-12}{12/10})
\end{soluzione}

\begin{soluzione}[Domanda~\ref{quiz:m27}]
\textbf{(d)}. Ogni istante deve essere un vettore di feature (qui di dimensione 1): \code{x.unsqueeze(-1)}. (rif.\ \rif{2026-10-12}{12/10})
\end{soluzione}

\begin{soluzione}[Domanda~\ref{quiz:m28}]
\textbf{(b)}. (rif.\ \rif{2026-10-14}{14/10})
\end{soluzione}

\begin{soluzione}[Domanda~\ref{quiz:m29}]
\textbf{(c)}. Riordinando i nodi cambia la matrice; un nodo in più richiede un neurone in più. (rif.\ \rif{2026-10-14}{14/10})
\end{soluzione}

\begin{soluzione}[Domanda~\ref{quiz:m30}]
\textbf{(a)}. (rif.\ \rif{2026-10-14}{14/10})
\end{soluzione}

\begin{soluzione}[Domanda~\ref{quiz:m31}]
\textbf{(d)}. (a) è GAT; (b) GCN. (rif.\ \rif{2026-10-14}{14/10})
\end{soluzione}

\begin{soluzione}[Domanda~\ref{quiz:m32}]
\textbf{(b)}. \code{batch.x.size(0)} conta i nodi (es.\ 288), non i grafi (2). (rif.\ \rif{2026-10-19}{19/10})
\end{soluzione}

\begin{soluzione}[Domanda~\ref{quiz:m33}]
\textbf{(c)}. (rif.\ \rif{2026-10-19}{19/10})
\end{soluzione}

\begin{soluzione}[Domanda~\ref{quiz:m34}]
\textbf{(a)}. È l'equivalente efficiente di one-hot + layer lineare; random forest non può imparare embedding. (rif.\ \rif{2026-10-19}{19/10})
\end{soluzione}

\begin{soluzione}[Domanda~\ref{quiz:m35}]
\textbf{(d)}. I sopravvissuti sono scalati di $1/(1-p) = 2$: $3\cdot 2 = 6$; il valore atteso resta 3. (rif.\ \rif{2026-10-20}{20/10})
\end{soluzione}

\begin{soluzione}[Domanda~\ref{quiz:m36}]
\textbf{(b)}. (rif.\ \rif{2026-10-20}{20/10})
\end{soluzione}

\begin{soluzione}[Domanda~\ref{quiz:m37}]
\textbf{(c)}. Il valore è normale, è il contesto (giorno festivo) a renderlo anomalo. (rif.\ \rif{2026-11-02}{02/11})
\end{soluzione}

\begin{soluzione}[Domanda~\ref{quiz:m38}]
\textbf{(a)}. (rif.\ \rif{2026-11-02}{02/11})
\end{soluzione}

\begin{soluzione}[Domanda~\ref{quiz:m39}]
\textbf{(d)}. L'SSE dovrebbe scendere con $K$: una risalita indica una cattiva inizializzazione. (rif.\ \rif{2026-11-02}{02/11})
\end{soluzione}

\begin{soluzione}[Domanda~\ref{quiz:m40}]
\textbf{(b)}. (rif.\ \rif{2026-11-02}{02/11})
\end{soluzione}

\begin{soluzione}[Domanda~\ref{quiz:m41}]
\textbf{(c)}. Varianza in byte (milioni) contro varianza in millisecondi: va applicato \code{StandardScaler} prima della PCA. (rif.\ \rif{2026-11-09}{09/11})
\end{soluzione}

\begin{soluzione}[Domanda~\ref{quiz:m42}]
\textbf{(a)}. Il modello si specializza sul normale solo se vede soltanto il normale; con anomalie in training le ricostruisce. (rif.\ \rif{2026-11-09}{09/11})
\end{soluzione}

\begin{soluzione}[Domanda~\ref{quiz:m43}]
\textbf{(d)}. (rif.\ \rif{2026-11-09}{09/11})
\end{soluzione}

\begin{soluzione}[Domanda~\ref{quiz:m44}]
\textbf{(b)}. Un attacco etichettato non è per forza un outlier di densità: servirebbe un $k$ più grande del gruppo. (rif.\ \rif{2026-11-09}{09/11})
\end{soluzione}

\begin{soluzione}[Domanda~\ref{quiz:m45}]
\textbf{(b)}. È la baseline usata nel Lab~2 sulle API call: solida, ma senza ordine né semantica. (rif.\ \rif{2026-11-23}{23/11})
\end{soluzione}

\begin{soluzione}[Domanda~\ref{quiz:m46}]
\textbf{(c)}. La (a) è il CBOW. (rif.\ \rif{2026-11-23}{23/11})
\end{soluzione}

\begin{soluzione}[Domanda~\ref{quiz:m47}]
\textbf{(b)}. È una softmax su $V$ classi; il negative sampling tiene il contributo positivo e usa $k \ll V$ negativi (più $k$ è grande, migliore l'approssimazione). (rif.\ \rif{2026-11-23}{23/11})
\end{soluzione}

\begin{soluzione}[Domanda~\ref{quiz:m48}]
\textbf{(c)}. Word2Vec è \textbf{non contestuale}: un vettore per parola. La (b) descrive un modello contestuale (transformer). (rif.\ \rif{2026-11-23}{23/11})
\end{soluzione}

\begin{soluzione}[Domanda~\ref{quiz:m49}]
\textbf{(c)}. La darknet non risponde, quindi vede solo IP e porte, senza label. (rif.\ \rif{2026-11-23}{23/11})
\end{soluzione}

\begin{soluzione}[Domanda~\ref{quiz:m50}]
\textbf{(b)}. Una frase per servizio, perché le porte contano: IP che sondano servizi diversi non vanno messi nello stesso contesto. (cap.~\ref{cap:darkvec}, slide 8)
\end{soluzione}

\begin{soluzione}[Domanda~\ref{quiz:m51}]
\textbf{(b)}. (cap.~\ref{cap:darkvec}, slide 8)
\end{soluzione}

\begin{soluzione}[Domanda~\ref{quiz:vfA} (gruppo A)]
\textbf{A1 F, A2 V, A3 F, A4 V}.\\
A1: l'epoch è il passaggio di tutto il dataset; con i mini-batch ci sono molti aggiornamenti per epoch. A3: loss 0 è sospetta (overfitting). (rif.\ \rif{2026-09-30}{30/09})
\end{soluzione}

\begin{soluzione}[Domanda~\ref{quiz:vfB} (gruppo B)]
\textbf{B1 V, B2 F, B3 V, B4 F}.\\
B2: i vettori one-hot sono equidistanti. B4: è la media \en{macro} a pesare ugualmente le classi. (rif.\ \rif{2026-10-05}{05/10})
\end{soluzione}

\begin{soluzione}[Domanda~\ref{quiz:vfC} (gruppo C)]
\textbf{C1 F, C2 V, C3 V, C4 F}.\\
C1: è \code{model.eval()} a spegnere il dropout. C4: si passano gli indici di classe interi (\code{torch.long}). (rif.\ \rif{2026-10-06}{06/10})
\end{soluzione}

\begin{soluzione}[Domanda~\ref{quiz:vfD} (gruppo D)]
\textbf{D1 F, D2 F, D3 V, D4 V}.\\
D1: i pesi sono condivisi tra i passi. D2: richiede la sequenza completa, il futuro non è disponibile in tempo reale. (rif.\ \rif{2026-10-12}{12/10}, \rif{2026-10-14}{14/10})
\end{soluzione}

\begin{soluzione}[Domanda~\ref{quiz:vfE} (gruppo E)]
\textbf{E1 V, E2 F, E3 F, E4 V}.\\
E2: le GNN sono induttive. E3: il vocabolario va costruito solo sul training (leakage); gli sconosciuti vanno in \code{<UNK>}. (rif.\ \rif{2026-10-14}{14/10}, \rif{2026-10-19}{19/10}, \rif{2026-10-20}{20/10})
\end{soluzione}

\begin{soluzione}[Domanda~\ref{quiz:vfF} (gruppo F)]
\textbf{F1 V, F2 F, F3 V, F4 F}.\\
F2: servono feature aggregate su finestre. F4: si addestra solo su dati normali, ma senza passare label. (rif.\ \rif{2026-11-02}{02/11})
\end{soluzione}

\begin{soluzione}[Domanda~\ref{quiz:vfG} (gruppo G)]
\textbf{G1 F, G2 F, G3 V, G4 F}.\\
G1: le label servono solo a validare (anomalie note). G2: t-SNE preserva i vicini, non le distanze. G4: con forte sbilanciamento la ROC-AUC può ingannare come l'accuracy: guarda precision, recall e F1. (rif.\ \rif{2026-11-09}{09/11})
\end{soluzione}

\begin{soluzione}[Domanda~\ref{quiz:vfH} (gruppo H)]
\textbf{H1 F, H2 V, H3 F, H4 V}.\\
H1: è self-supervised, la supervisione viene dal contesto delle parole. H3: è non contestuale; contestuali saranno i transformer. (rif.\ \rif{2026-11-23}{23/11})
\end{soluzione}
